\scanpage{70}\mspage{69}
(ce qui est \unsure{automatique} si on suppose $\unsure{\widehat F|\widetilde F}$ fini),
alors je dis que le \deleted{\unsure{$F'_X$}} $\mathrm{Sup}$ $F'$ répond
à la question. En effet, on aura \unsure{$\Gamma$|$F$}
\[
\widetilde F'=\deleted{\unsure{S}}\bigcup_X\widetilde F'_X\quad,
\]
et on aura $F'\ll\unsure{F|\mathbb F|\overline F}$ par \unsure{At|Ax} 4. Pour
prouver que $F'\mid X=F'_X$\inserted{$\forall X\in\widetilde F$},
\deleted{\unsure{on note que}}
\[
\deleted{F'_X=\mathop{\mathrm{Sup}}\limits_{Z\in\unsure{F|\widetilde F}}\bigl(F'_Z\mid X\bigr)}
\]
% slanted
\margin{\deleted{\unsure{réduction} des raffinements \inserted{$F'_Z$} de $F$ sur \unsure{sous}-figure $X$ de $F$}}
\deleted{Et par la condition \unsure{mise|mises} sur le \unsure{syst.} des $F'_T$, on aura}
\diagram{\begin{tikzcd}[ampersand replacement=\&,column sep=2pt,row sep=1pt,cells={nodes={font=\scriptsize,inner sep=1pt}}]
{} \& {} \& {} \& {} \& {} \& {} \& \deleted{F'} \& \deleted{\ll} \& \deleted{F} \& {} \& {}\\
{} \& {} \& {} \& {} \& {} \& {} \& \deleted{\rotatebox{90}{$\leq$}} \& \deleted{\unsure{F'_Z|F_Z}} \& \deleted{\rotatebox{80}{$\leq$}} \& \deleted{\rotatebox{105}{$\leq$}} \& \deleted{F'_Z\mid X}\\
\deleted{\mathop{\mathrm{Sup}}\limits_Z F'_Z} \arrow[dd,dash,bend right=16] \& \deleted{=} \arrow[dd,dash,bend left=18] \& \deleted{F'} \& \deleted{\ll} \& \deleted{F} \& {} \& \deleted{F'_Z} \& \deleted{\ll} \& \deleted{Z} \& \deleted{X} \& {}\\
\deleted{\unsure{\rotatebox{90}{$\leq$}}} \& {} \& \deleted{\rotatebox{90}{$\leq$}} \& {} \& \deleted{\rotatebox{90}{$\leq$}} \& {} \& {} \& {} \& \deleted{\rotatebox{125}{$\leq$}} \& \deleted{\rotatebox{70}{$\leq$}} \& {}\\
\deleted{F'_Z\mid X} \& \deleted{\leq} \& \deleted{F'\unsure{\mathord{.}|\cdot}\mid X} \& \deleted{\ll} \& \deleted{X} \& {} \& {} \& {} \& {} \& |[alias=inkIntersection]| \deleted{X\cap Z} \& {}
\arrow[from=1-5,to=3-5,dash,start anchor=north,end anchor=center]
\arrow[from=1-8,to=3-9,dash,start anchor=north,end anchor=south west,bend right=7]
\arrow[from=1-10,to=3-10,dash,start anchor=north east,end anchor=east,bend right=5]
\arrow[from=3-3,to=5-1,dash,bend right=12]
\arrow[from=3-4,to=5-5,dash,bend right=25]
\arrow[from=3-5,to=2-11,dash,start anchor=center,end anchor=east,
  to path={.. controls +(0,-8mm) and +(-9mm,4mm) .. (inkIntersection.south west)
  .. controls +(6mm,-2mm) and +(2mm,5mm) .. (inkIntersection.east)
  .. controls +(-2mm,-2mm) and +(-4mm,-7mm) .. (\tikztotarget) \tikztonodes}]
\end{tikzcd}}

On regarde la \unsure{situation ensembliste}
\diagram{\begin{tikzcd}[ampersand replacement=\&,row sep=0pt,column sep=large]
\widetilde F'=\bigcup\widetilde F'_X \arrow[r] \& \widetilde F\,\unsure{-}\\
{} \& \rotatebox{90}{$\subset$}\\
{} \& \widetilde X
\end{tikzcd}}
où $\widetilde F$ est réunion des $\widetilde X$, et où $\widetilde F'$ a été
\unsure{obtenu} par recollement \deleted{\unsure{des|du}} \inserted{du syst. des} \unsure{ens} $\widetilde F'_X$
sur la famille admissible des $\widetilde X$ \unsure{recouvrant}
$\widetilde F$ ...

\uline{Corollaire} \uline{2}\quad Sous les conditions du Cor. 3
de la page 66, si on se donne un
\uline{\unsure{syst} compatible} $(F'_i)$ de raffinements des
$F_i$ (définir ce terme), alors
